\section*{Chapter \ref{ch:m2:corporate-bonds} --- Corporate Bonds}

\begin{solution}{exo:m2:corporate-bonds:1}
\omterm{def:m2:corporate-bonds:spreads}{G-spread} $6.10 - 4.68 = 142$ basis points; \omterm{def:m2:corporate-bonds:spreads}{I-spread} $6.10 - 4.40 = 170$.
\end{solution}

\begin{solution}{exo:m2:corporate-bonds:2}
It has a split rating: \omterm{def:m2:corporate-bonds:rating}{investment grade} at one agency, \omterm{def:m2:corporate-bonds:rating}{high yield} at the other. Whether
it counts as \omterm{def:m2:corporate-bonds:rating}{investment grade} depends on the rule of the index or mandate (some use
the lowest rating, some the middle of three). By the ECB's definition its issuer is a
\omterm{def:m2:corporate-bonds:rating}{fallen angel}, downgraded to \omterm{def:m2:corporate-bonds:rating}{high yield} by at least one major agency.
\end{solution}

\begin{solution}{exo:m2:corporate-bonds:3}
$-6 \times 0.0080 \times 100 = -4.8$ per 100.
\end{solution}

\begin{solution}{exo:m2:corporate-bonds:4}
Both measure the gap between the bond's price and its value on the swap curve, one as
a running spread on par, the other as a shift of every discount rate; near par and
for moderate spreads they agree to first order, and the 2.2 basis points are
compounding and weighting differences. They part as the spread grows and the price
moves away from par: repriced at 85, the bond has a \omterm{def:m2:corporate-bonds:spreads}{Z-spread} of 401.6 basis points and
an \omterm{def:m2:corporate-bonds:spreads}{asset-swap spread} of 375.5.
\end{solution}

\begin{solution}{exo:m2:corporate-bonds:5}
$0.0140/0.60 = 2.33\%$ a year. The rest of the spread, whatever actual defaults fall
short of that, pays for bearing the uncertainty of default and loss, for
illiquidity, and in some markets for taxes and for the cost of holding the bond.
\end{solution}

\begin{solution}{exo:m2:corporate-bonds:6}
Investors must absorb a large new supply at once and commit before the book is
complete; the issuer pays a little to fill it quickly and at size. If the concession
disappears as the bonds trade up to the company's curve, the buyers at issue gain it
and the issuer has paid it.
\end{solution}

\begin{solution}{exo:m2:corporate-bonds:7}
At 105 the yield is 4.6548\%: \omterm{def:m2:corporate-bonds:spreads}{G-spread} 15.5 basis points, \omterm{def:m2:corporate-bonds:spreads}{I-spread} 40.5, \omterm{def:m2:corporate-bonds:spreads}{Z-spread} 41.2,
\omterm{def:m2:corporate-bonds:spreads}{asset-swap spread} 43.2.
\end{solution}

\begin{solution}{exo:m2:corporate-bonds:8}
Part of the callable's spread pays for the call its holder has sold the issuer: its
\omterm{def:m2:corporate-bonds:oas}{option-adjusted spread}, the credit part, is lower. In
\cref{ex:m2:corporate-bonds:four}'s terms a \omterm{def:m2:corporate-bonds:spreads}{Z-spread} of 141 basis points became an OAS
of about 61 once a three-year call was valued. Compare OAS with OAS, not \omterm{def:m2:corporate-bonds:spreads}{Z-spread} with
\omterm{def:m2:corporate-bonds:spreads}{Z-spread}.
\end{solution}

\begin{solution}{pb:m2:corporate-bonds:1}
\textbf{1.} $6 \times 1.50\% = 9.0$ per 100.
\textbf{2.} $6 \times 0.80\% = 4.8$ per 100, 53\% of the fall.
\textbf{3.} Investors read the same results, leverage and outlook changes the
agencies do, and act on them continuously; the agencies act by committee and
announcement.
\textbf{4.} Sold early, or reduced its holding, if its mandate or risk policy treats a
negative outlook at BBB$-$ as a warning.
\textbf{5.} $6 \times 0.70\% = 4.2$ per 100.
\textbf{6.} 50 basis points over the fair 350: 3.0 per 100.
\textbf{7.} About USD~1.5 million on USD~50 million, sold at 3 points below fair value.
\textbf{8.} High-yield funds and investors who can hold \omterm{def:m2:corporate-bonds:rating}{high yield}; they must take on
a large supply at once and fund it by selling other bonds, and they know the sellers
are forced.
\textbf{9.} Index providers' rules differ, so not every index drops the bond on the
same day; funds that sample the index and may keep a bond for a while spread their
sales out, reducing the pressure.
\textbf{10.} The overshoot, if the spread returns to 350: 3 points, USD~1.5 million,
plus the higher coupon income.
\textbf{11.} 3.0 per 100, from 400 to 350.
\textbf{12.} Further downgrades or default, and a spread that does not recover.
\textbf{13.} That most \omterm{def:m2:corporate-bonds:rating}{fallen angels} survive (nearly a quarter returned to \omterm{def:m2:corporate-bonds:rating}{investment
grade}, almost half stayed \omterm{def:m2:corporate-bonds:rating}{high yield}) but 12\% defaulted within the twenty years.
\textbf{14.} About $4.00\% \times 0.25 = 1.0$ per 100 over three months in spread
carry, on top of the price gain.
\textbf{15.} The buyer is paid for taking the risk the sellers must shed, and the
spread may widen further before it recovers.
\textbf{16.} Because investment mandates, indices and regulations are written in
ratings: the line decides who may own the bond, whatever its price.
\textbf{17.} By offering a buyer for recently downgraded bonds within limits, it
would reduce the overshoot, as the 2020 facility's BB$-$/Ba3 provision was designed to.
\textbf{18.} To keep a \omterm{def:m2:corporate-bonds:rating}{fallen angel} for a period, or up to a limit, instead of selling
at once.
\textbf{19.} Named result: \emph{the \omterm{def:m2:corporate-bonds:rating}{fallen angel}'s discount}: forced selling pushed
the bonds 50 basis points past fair value, 3 points per 100, USD~1.5 million on the
holding, which a buyer recovered within three months.
\textbf{20.} Because many holders must sell at the same moment to buyers who need a
discount to absorb them.
\end{solution}

\begin{solution}{iq:m2:corporate-bonds:1}
The \omterm{def:m2:corporate-bonds:spreads}{G-spread} is the yield over the interpolated government yield; the \omterm{def:m2:corporate-bonds:spreads}{Z-spread} is the
constant spread over the zero curve that reprices all the cash flows; the \omterm{def:m2:corporate-bonds:spreads}{asset-swap
spread} is what an investor earns over the floating rate by swapping the bond's
coupons at par. They differ by the \omterm{def:m2:interest-rate-swaps:spread}{swap spread}, the curve's shape and the bond's
price away from par.
\iqlookfor{definitions and why they differ.}
\end{solution}

\begin{solution}{iq:m2:corporate-bonds:2}
An issuer downgraded from \omterm{def:m2:corporate-bonds:rating}{investment grade} to \omterm{def:m2:corporate-bonds:rating}{high yield}. Its bonds leave
investment-grade indices and many mandates, so holders must sell to a smaller set of
buyers who want a discount; the market also reprices the credit before the downgrade.
Evidence shows much of the move before the event and a partial recovery after.
\iqlookfor{the forced-selling mechanism and the timing.}
\end{solution}

\begin{solution}{iq:m2:corporate-bonds:3}
Estimate expected loss from default probabilities and recoveries, historical or
from a model, by rating and maturity; subtract it from the spread over a risk-free
curve; the remainder is compensation for risk and illiquidity, which can be related
to market volatility, liquidity measures and investors' capacity. Check the result
across ratings and through the cycle.
\iqlookfor{expected loss versus premium, and its time variation.}
\end{solution}

\begin{solution}{iq:m2:corporate-bonds:4}
Start from the issuer's other bonds and CDS, interpolated to the maturity, and from
comparable issuers; adjust for size, coupon, \omterm{def:m2:corporate-bonds:covenant}{covenants} and liquidity; check recent
trades and quotes on platforms; then price the risk of holding USD~50 million until
you can hedge or sell it, which widens the price with size and illiquidity.
\iqlookfor{relative value from the curve and comparables, plus a size charge.}
\end{solution}

\begin{solution}{iq:m2:corporate-bonds:5}
The \omterm{def:m2:corporate-bonds:spreads}{Z-spread} includes payment for the call the holder has sold; the OAS removes the
option's value, computed with a rate model, and keeps the credit and liquidity part.
Compare OAS across bonds with different options; \omterm{def:m2:corporate-bonds:spreads}{Z-spreads} only across bullets.
\iqlookfor{option value inside the spread.}
\end{solution}

\begin{solution}{iq:m2:corporate-bonds:6}
Maintain curves (Treasury, swap, zero) updated on market events; hold each bond's
cash-flow schedule precomputed; recompute prices and spreads in vectorised batches,
with root finders warm-started from the last value; parallelise across bonds and
recompute only those whose price or curve changed; validate against a reference
implementation and flag failures to converge.
\iqlookfor{precomputation, warm starts, incremental updates.}
\end{solution}
